\section{The radius-weighted block estimate}\label{s:weighted-proof}

We now develop the cap-moment method for arbitrary bounded continuation
functions. Weighting the mean on an affine subsphere by a power of its
radius makes the same test useful at every scale. The resulting block
inequality is independent of a learner and also gives a precision bound
at success probability one half.

For a Borel function \(g:S^{d-1}\to[0,1]\), an exponent \(a=d/16\), and an integer
\(1\le n\le d\), define
\begin{equation}\label{s:weighted-norm}
 B_n(g)=\sup_{\substack{\dim H=n\\ z\perp H,\ R>0\\
                         \norm z^2+R^2=1}}
 R^a\int g(z+R\omega)\,d\sigma_H(\omega).
\end{equation}
In particular \(B_d(g)=\int g\,d\sigma_{\R^d}\). The radius weight
allows small subspheres to carry more mass while keeping the same test
at every scale.

The block may choose among arbitrary bounded continuation functions.

\begin{theorem}[Radius-weighted block estimate]\label{s:weighted-theorem}
There are absolute constants \(C\ge1\) and \(d_*\) such that the
following holds for every integer \(d\ge d_*\). Put
\[
 k=m=\lfloor d/16\rfloor,\qquad a=d/16.
\]
Let \(W\ge1\), let \(\mathcal V\) be a nonempty finite set with
\(|\mathcal V|\le W\), and let
\(\kappa_v:\R^{k\times d}\times\R^k\to[0,1]\) be Borel functions
with \(\sum_{v\in\mathcal V}\kappa_v\le1\). For Borel
\(g_v:S^{d-1}\to[0,1]\), set
\[
 g_{\rm in}(s)=
 \E_X\sum_{v\in\mathcal V}\kappa_v(X,Xs)g_v(s),
\]
where \(X\) has \(k\) independent \(N(0,I_d)\) rows. For every integer
\(n\) with \(d/2\le n\le d\),
\begin{equation}\label{s:weighted-block}
 B_n(g_{\rm in})
 \le W^{1/m}C^d\max_{v\in\mathcal V}B_{n-2}(g_v).
\end{equation}
If \(\epsilon>0\), \(\pi\) is any probability measure on \(S^{d-1}\) and
\[
 g_\pi(s)=\int\ind_{\{\arccos\ip us\le\epsilon\}}\,\pi(du),
\]
then, in the same range of \(n\),
\begin{equation}\label{s:weighted-terminal}
 B_n(g_\pi)\le C^d\epsilon^a.
\end{equation}
\end{theorem}

The width factor pays for selecting a continuation after the block.
The two-dimensional change instead records the latitude tests used to
bound each conditional direction average. We first establish the moment
estimate that connects these two features.

\subsection{Caps and subspace moments}

The use of replica moments and successive orthogonalization is related
to Sharan, Sidford, and Valiant \cite[Section~7, Lemma~12]{SSV2019}.
The density becomes large near the preceding span. The next lemma
shows that a cap estimate with sufficient dimension slack makes this
singularity integrable. Its explicit constant retains a useful
quantitative form of the argument.

\begin{lemma}[Caps and subspace moments]\label{s:cap-moment}
Let \(E\) have dimension \(N\le d\), and let \(L\) be a uniform
\(K\)-space in \(E\), with integers \(1\le h<K\le N\). Put
\(b=N-K\). Let \(f:S(E)\to[0,1]\) be Borel. Suppose that for some
finite \(\Lambda\ge0\), constants \(A_0,A_1\ge1\), and
\(0\le a<N-1\),
\begin{align}
 \int f\,d\sigma_E&\le A_0^d\Lambda,\label{s:cap-total}\\
 \int_{\norm{y-y_0}\le\delta}f(y)\,d\sigma_E(y)
 &\le A_1^d\Lambda\delta^{N-1-a}
 \quad(y_0\in S(E),\ \delta>0).\label{s:cap-hypothesis}
\end{align}
If
\begin{equation}\label{s:moment-margin}
 N-1-a-(h-1)-b\ge1,
\end{equation}
then
\begin{equation}\label{s:moment-bound}
 \E_L\left(\int f\,d\sigma_L\right)^h
 \le\left([4\max\{A_0,9A_1\}]^d\Lambda\right)^h.
\end{equation}
For a nonempty family of at most \(W\) functions satisfying the same
conditions with numbers \(\Lambda_v\),
\begin{equation}\label{s:moment-maximum}
 \E_L\max_v\int f_v\,d\sigma_L
 \le W^{1/h}[4\max\{A_0,9A_1\}]^d\max_v\Lambda_v.
\end{equation}
In particular, \(A_0=1,A_1=4\) gives the constant \(144^d\).
\end{lemma}
\begin{proof}
Put \(\alpha=N-1-a\) and \(Q_0=\max\{A_0,9A_1\}\).
For a \(j\)-dimensional subspace \(V\subset E\), \(1\le j<h\),
and \(0<\delta<1\), every sphere point within distance \(\delta\)
of \(V\) lies within \(2\delta\) of its normalized projection onto
\(S(V)\). A maximal \(\delta\)-separated subset of \(S(V)\) has at
most \((3/\delta)^j\) points: the open balls of radius \(\delta/2\)
about them are disjoint inside the ball of radius \(1+\delta/2\)
in \(V\). Caps of radius \(3\delta\) about these points cover the
tube. Hence
\begin{equation}\label{s:tube}
 \int_{\dist(y,V)\le\delta}f\,d\sigma_E
 \le A_1^d3^{\alpha+j}\Lambda\delta^{\alpha-j}
 \le Q_0^d\Lambda\delta^{\alpha-j}.
\end{equation}
The last inequality uses \(\alpha,j\le d\). At \(\delta=1\), the
same bound follows from \eqref{s:cap-total}. The zero-distance set
has sphere measure zero because \(j<N\). The shells
\(2^{-i-1}<\dist(y,V)\le2^{-i}\) therefore give
\begin{align}
 \int f(y)\dist(y,V)^{-b}\,d\sigma_E(y)
 &\le Q_0^d\Lambda\,2^b
       \sum_{i\ge0}2^{-i(\alpha-j-b)}\notag\\
 &\le Q_0^d\Lambda\,2^{b+1}
 \le (4Q_0)^d\Lambda.\label{s:inverse-distance}
\end{align}
Here \(\alpha-j-b\ge1\) follows from
\eqref{s:moment-margin}, and \(b+1\le2d\). For \(j=0\), the distance
of a unit vector from \(\{0\}\) is one and the same conclusion follows
directly from total mass. Thus this case needs neither a net nor a
shell sum.

Conditional on \(L\), draw \(Y_1,\ldots,Y_h\) independently from
\(\sigma_L\). Expanding the power in \eqref{s:moment-bound} gives
\(\E\prod_{i=1}^h f(Y_i)\).
Lemma~\ref{s:directions} and \eqref{s:inverse-distance} bound the last
conditional factor by \((4Q_0)^d\Lambda\). Continue in reverse
conditional order, using total mass for the first direction. This
proves \eqref{s:moment-bound}. Finally, writing
\(F_v(L)=\int f_v\,d\sigma_L\), Jensen's inequality gives
\[
 \E\max_vF_v(L)
 \le\E\left(\sum_vF_v(L)^h\right)^{1/h}
 \le\left(\sum_v\E F_v(L)^h\right)^{1/h},
\]
which proves \eqref{s:moment-maximum}.
\end{proof}

The lemma pays for the span of prior replicas through the tube
exponent. It does not spend a new induction dimension for each
replica: every conditional factor is tested against the same cap
hypothesis on \(S(E)\). In the applications, one latitude of an
affine subsphere supplies that hypothesis.

\subsection{Proof of the block estimate and its application}

A latitude in the perpendicular ambient sphere has two fewer linear
dimensions than the original test. Its radius-weighted mean supplies the cap hypothesis;
a final beta integral averages the residual radius.

\begin{proof}[Proof of Theorem~\ref{s:weighted-theorem}]
We work at sufficiently large \(d\), with the parameters in the
theorem. First consider a terminal function \(g_\pi\) and a test
\((z,H,R)\) of dimension \(n\ge d/2\). If \(R\ge2\epsilon\), the
successful directions for any fixed output lie in a chordal cap of
radius \(2\epsilon/R\le1\). By \eqref{s:sharp-cap}, their mass is at
most \((c_n/2)(2\epsilon/R)^{n-1}\). Since \(n-1\ge a\),
multiplication by \(R^a\) bounds this by
\((c_n/2)(2\epsilon)^a\le C^d\epsilon^a\).
If \(R<2\epsilon\), total mass gives \(R^a\le(2\epsilon)^a\).
Averaging the output law proves \eqref{s:weighted-terminal}.

For the block bound, fix a test \((z,H,R)\) of dimension \(n\).
Let \(V\) be the span in \(H\) of the \(k\) projected rows. Put
\[
 p=R P_V\omega,\qquad r=\sqrt{R^2-\norm p^2},\qquad
 K=H\cap V^\perp,\qquad q=n-k.
\]
Almost surely \(p\ne0\) and \(r>0\). Conditional on the matrix and
\(p\), the target is \(z+p+rU\) with \(U\sim\sigma_K\), and its
labels are \(X(z+p)\). Lemma~\ref{s:state-maximum} bounds the
averaged block function by
\begin{equation}\label{s:weighted-max}
 \E_p\E_{K\mid p}\max_v
       \int g_v(z+p+rU)\,d\sigma_K(U).
\end{equation}
Here \(K\mid p\) is uniform among the \(q\)-spaces of
\(H'=H\cap p^\perp\), whose dimension is \(N=n-1\).
Indeed, the length of \(p\) fixes the beta fraction, its direction
fixes \(e\), and Lemma~\ref{s:residual} gives this conditional law.

Fix \(p\) and one function \(g=g_v\). On \(S(H')\) set
\[
 f(U)=g(z+p+rU),\qquad \Lambda=r^{-a}B_{n-2}(g).
\]
For \(w\in S(H')\), condition on the latitude coordinate
\(t=\ip Uw\in(-1,1)\). The target on that latitude is uniform on
\[
 z+p+rtw+r\sqrt{1-t^2}\,S(H'\cap w^\perp).
\]
This is an admissible affine subsphere of linear dimension \(n-2\):
its center is perpendicular to \(H'\cap w^\perp\), and its squared
center norm plus squared radius is
\(\norm z^2+\norm p^2+r^2=1\). The definition of \(B_{n-2}\) therefore
bounds the latitude average of \(f\) by
\(\Lambda(1-t^2)^{-a/2}\).

The coordinate density in dimension \(N\) has remaining exponent
\((N-3-a)/2\ge0\). Integrating it over the whole interval and then
over the cap interval gives
\begin{align}
 \int f\,d\sigma_{H'}&\le2c_N\Lambda,\label{s:weighted-total}\\
 \int_{\norm{U-w}\le\eta} f(U)\,d\sigma_{H'}(U)
 &\le c_N\Lambda\eta^{N-1-a}
 \qquad(0<\eta\le1).\label{s:weighted-cap}
\end{align}
For the second line the coordinate interval has length at most
\(\eta^2/2\) and \(1-t^2\le\eta^2\) there. For \(\eta>1\), total
mass extends the cap bound with an absolute exponential factor.
Since \(2c_N\le2\sqrt d\le4^d\), these estimates satisfy
Lemma~\ref{s:cap-moment} with absolute \(A_0,A_1\).

Use that lemma with ambient dimension \(N=n-1\), subspace dimension
\(q=n-k\), and moment \(m=k\). The singular power is
\(N-q=k-1\). The required inequalities hold at large \(d\):
\(1\le m<q\le N\), \(0\le a<N-1\), and, for
\(0\le l<m\),
\[
 N-1-a-l-(N-q)
 =n-1-a-l-k
 \ge n-a-2k
 \ge5d/16\ge1.
\]
The moment conclusion is
\begin{equation}\label{s:weighted-moment}
 \E_{K\mid p}\left(\int g(z+p+rU)\,d\sigma_K(U)\right)^m
 \le\left(C^d r^{-a}B_{n-2}(g)\right)^m.
\end{equation}
Each replica uses the same \(n-2\)-dimensional latitude tests.
No additional induction dimension is lost while the replicas are
integrated.

The maximum inequality \eqref{s:moment-maximum} now bounds the inner
expectation in \eqref{s:weighted-max} by
\[
 W^{1/m}C^d r^{-a}\max_v B_{n-2}(g_v).
\]
It remains to average the residual radius. By
Lemma~\ref{s:residual}, \(t=(r/R)^2\sim\Beta(q/2,k/2)\).
The denominator of its beta integral satisfies
\[
 \mathrm B(q/2,k/2)
 \ge \frac12\,4^{-(q/2-1)_+-(k/2-1)_+}
 \ge 4^{-d},
\]
by restriction to \([1/4,3/4]\), where \(x_+=\max\{x,0\}\).
After multiplication by \(t^{-a/2}\), the numerator is
\begin{equation}\label{s:weighted-beta-integral}
 \int_0^1 t^{(q-a)/2-1}(1-t)^{k/2-1}\,dt\le2.
\end{equation}
Indeed \(q-a\ge2\) makes the first exponent nonnegative. If \(k\ge2\)
the second is also nonnegative; if \(k=1\), its integral is at most
\(\int_0^1(1-t)^{-1/2}\,dt=2\). Thus
\(\E(R/r)^a\le C^d\). This direct beta-integral argument does not
require \(a/2\) to be an integer. Multiplying
\eqref{s:weighted-max} by \(R^a\), averaging, and taking the supremum
over the original test proves \eqref{s:weighted-block}.
\end{proof}

The block theorem now gives a complete lower bound from these norms.
This application uses success \(1/2\), which is enough for the
constant-accuracy estimate \eqref{s:quarter-endpoint}.

\begin{corollary}[Precision bound from radius weights]\label{s:weighted-application}
There is an absolute \(c>0\) such that, for every \(M(d)=o(d^2)\),
all sufficiently large \(d\), and all \(0<\epsilon\le1/10\), a learner
with uniform-prior angular success at least \(1/2\) satisfies
\[
 T\ge c\,d\log_2(1/\epsilon).
\]
The eventual threshold may depend on the memory sequence.
\end{corollary}
\begin{proof}
Put \(L=\log_2(1/\epsilon)\). The dimensions with \(T\ge dL\)
already satisfy the conclusion. On the remaining dimensions,
Lemma~\ref{s:short-run}, using
\eqref{s:sharp-short-run}, gives \(L=o(d)\), \(T=o(d^2)\), and
\(\log_2 W_*=o(d^2)\). Pad to \(b=\lceil T/k\rceil=o(d)\) blocks
of \(k=\lfloor d/16\rfloor\) samples with width \(W_*\), as in
Lemma~\ref{s:terminal-count}. A short final block is filled with
unused fresh rows.

Let \(G_{i,v}\) be the Borel suffix success from state \(v\) after
\(i\) blocks. Recursion \eqref{s:recursion} and
Theorem~\ref{s:weighted-theorem} apply at each boundary. Use dimension
\(n=d-2i\) at stage \(i\). For all sufficiently large dimensions
under consideration, every needed dimension, including the terminal
one, satisfies
\begin{equation}\label{s:weighted-budget}
 d-2i\ge d-2b\ge d/2\qquad(0\le i\le b).
\end{equation}
Iteration of \eqref{s:weighted-block}, followed by
\eqref{s:weighted-terminal}, gives
\begin{equation}\label{s:weighted-iteration}
 \frac12\le
       [W_*^{1/m}C^d]^b C^d\epsilon^a.
\end{equation}
The same bound holds for each initial state and each fixed independent
choice of rules, so averaging them is legitimate.

At large \(d\), \(m\ge d/32\) and
\(\log_2 W_*/m\le d\). Taking base-two logarithms in
\eqref{s:weighted-iteration}, with \(a=d/16\), gives
\(L\le C_1(b+1)\) for an absolute \(C_1\ge1\).
If \(L\ge4C_1\), then \(b\le T/k+1\) implies
\[
 T\ge kL/(2C_1).
\]
Since \(k\ge d/32\) eventually, this is the asserted rate. For
\(L<4C_1\), \eqref{s:quarter-endpoint} gives the same rate with
an absolute smaller constant. This completes every accuracy regime.
\end{proof}
